\documentclass[../../main.tex]{subfiles}% 注意这里的文件路径不能用 ./main.tex ,否则用latexmk编译子文件会报错
\graphicspath{{\subfix{./image/}}{\subfix{../../image/}}} % 指定图片引用路径,后续可以直接使用图片文件名
% 如果图片存放在上级目录,则使用 ../../image/ 作为备用路径

% 例如:
% \begin{figure}[H]
% \centering
% \includegraphics[scale=0.3]{图.png}
% \caption{}
% \label{figure:图}
% \end{figure}
% 注意:上述\label{}一定要放在\caption{}之后,否则引用图片序号会只会显示??.

\begin{document}

\section{曲面的唯一性定理}

\begin{theorem}
设$S_1,S_2$是定义在同一参数区域$D \subset E^2$上的两个正则参数曲面。若在每一点$(u^1,u^2) \in D$，曲面$S_1$和$S_2$都有相同的第一基本形式和第二基本形式，则曲面$S_1$和$S_2$在空间$E^3$的一个刚体运动下是彼此重合的。
\end{theorem}

\begin{remark}
该定理称为\textbf{曲面的唯一性定理}，说明曲面不计空间位置时，由其第一基本形式和第二基本形式唯一确定，具有重要的理论意义。
\end{remark}

\begin{proof}

$S_1,S_2$采用同一组参数，二者具有相同的第一、第二基本形式等价于在每一点具有相同的第一类、第二类基本量。

取定点$u_0=(u_0^1,u_0^2) \in D$，设曲面$S_i$的自然标架场为$\{\boldsymbol{r}^{(i)};\boldsymbol{r}_1^{(i)},\boldsymbol{r}_2^{(i)},\boldsymbol{n}^{(i)}\}\ (i=1,2)$，在$u_0$处满足
\begin{align}
\begin{split}
\boldsymbol{r}_\alpha^{(1)}(u_0) \cdot \boldsymbol{r}_\beta^{(1)}(u_0) &= \boldsymbol{r}_\alpha^{(2)}(u_0) \cdot \boldsymbol{r}_\beta^{(2)}(u_0) = g_{\alpha\beta}(u_0),\\
\boldsymbol{r}_\alpha^{(1)}(u_0) \cdot \boldsymbol{n}^{(1)}(u_0) &= \boldsymbol{r}_\alpha^{(2)}(u_0) \cdot \boldsymbol{n}^{(2)}(u_0) = 0,\\
\boldsymbol{n}^{(1)}(u_0) \cdot \boldsymbol{n}^{(1)}(u_0) &= \boldsymbol{n}^{(2)}(u_0) \cdot \boldsymbol{n}^{(2)}(u_0) = 1,
\end{split}
\label{eq:::IOJIOWAIJ2.1}
\end{align}
且两标架均为右手系。故存在$E^3$中的刚体运动$\sigma$，将标架$\{\boldsymbol{r}^{(2)}(u_0);\boldsymbol{r}_1^{(2)}(u_0),\boldsymbol{r}_2^{(2)}(u_0),\boldsymbol{n}^{(2)}(u_0)\}$映为$\{\boldsymbol{r}^{(1)}(u_0);\boldsymbol{r}_1^{(1)}(u_0),\boldsymbol{r}_2^{(1)}(u_0),\boldsymbol{n}^{(1)}(u_0)\}$。

刚体运动保持第一类、第二类基本量不变，故$\sigma(S_2)$与$S_1$处处具有相同的基本量，二者自然标架场均满足同一组偏微分方程。不妨记$\sigma(S_2)$仍为$S_2$，则$S_1,S_2$在$u_0$处自然标架重合，且处处基本量相同。

定义函数
\begin{align}
\begin{split}
f_{\alpha\beta}(u) &= (\boldsymbol{r}_\alpha^{(1)} - \boldsymbol{r}_\alpha^{(2)}) \cdot (\boldsymbol{r}_\beta^{(1)} - \boldsymbol{r}_\beta^{(2)}),\\
f_\alpha(u) &= (\boldsymbol{r}_\alpha^{(1)} - \boldsymbol{r}_\alpha^{(2)}) \cdot (\boldsymbol{n}^{(1)} - \boldsymbol{n}^{(2)}),\\
f(u) &= (\boldsymbol{n}^{(1)} - \boldsymbol{n}^{(2)})^2,
\end{split}
\label{eq:::IOJIOWAIJ2.2}
\end{align}
由$u_0$处标架重合，可得初值条件
\begin{align}
f_{\alpha\beta}(u_0) = 0,\quad f_\alpha(u_0) = 0,\quad f(u_0) = 0.
\label{eq:::IOJIOWAIJ2.3}
\end{align}

结合自然标架的运动公式，可推得$f_{\alpha\beta},f_\alpha,f$满足一阶线性齐次偏微分方程组
\begin{align}
\begin{cases}
\displaystyle\frac{\partial f_{\alpha\beta}}{\partial u^\gamma} = \Gamma_{\gamma\alpha}^\delta f_{\delta\beta} + \Gamma_{\gamma\beta}^\delta f_{\delta\alpha} + b_{\gamma\alpha}f_\beta + b_{\gamma\beta}f_\alpha,\\[6pt]
\displaystyle\frac{\partial f_\alpha}{\partial u^\gamma} = -b_\gamma^\delta f_{\delta\alpha} + \Gamma_{\gamma\alpha}^\delta f_\delta + b_{\gamma\alpha}f,\\[6pt]
\displaystyle\frac{\partial f}{\partial u^\gamma} = -2b_\gamma^\delta f_\delta.
\end{cases}
\label{eq:::IOJIOWAIJ2.4}
\end{align}

一阶线性齐次偏微分方程组在给定初值条件下解唯一，零解是其解，故
\begin{align}
f_{\alpha\beta}(u) = 0,\quad f_\alpha(u) = 0,\quad f(u) = 0.
\label{eq:::IOJIOWAIJ2.5}
\end{align}

由此可得
$$|\boldsymbol{r}_\alpha^{(1)}(u) - \boldsymbol{r}_\alpha^{(2)}(u)|^2 = f_{\alpha\alpha}(u) = 0,\quad |\boldsymbol{n}^{(1)}(u) - \boldsymbol{n}^{(2)}(u)|^2 = f(u) = 0,$$
即
\begin{align}
\boldsymbol{r}_\alpha^{(1)}(u) = \boldsymbol{r}_\alpha^{(2)}(u),\ \alpha=1,2;\quad \boldsymbol{n}^{(1)}(u) = \boldsymbol{n}^{(2)}(u).
\label{eq:::IOJIOWAIJ2.6}
\end{align}

再定义
\begin{align}
h(u) = (\boldsymbol{r}^{(1)}(u) - \boldsymbol{r}^{(2)}(u))^2,
\label{eq:::IOJIOWAIJ2.7}
\end{align}
对$u^\gamma$求导并结合\eqref{eq:::IOJIOWAIJ2.6}得
$$\frac{\partial h}{\partial u^\gamma} = 2(\boldsymbol{r}^{(1)}(u) - \boldsymbol{r}^{(2)}(u)) \cdot (\boldsymbol{r}_\gamma^{(1)}(u) - \boldsymbol{r}_\gamma^{(2)}(u)) = 0,$$
故$h(u)=h(u_0)=0$，即$\boldsymbol{r}^{(1)}(u)=\boldsymbol{r}^{(2)}(u)$，因此曲面$S_1$与$S_2$重合。

\end{proof}

\begin{theorem}
设$S_i\ (i=1,2)$是空间$E^3$中的两个正则参数曲面，其第一基本形式和第二基本形式分别为$\mathrm{I}_i$和$\mathrm{II}_i$。若存在光滑映射$\sigma: S_1 \to S_2$，使得
\begin{align}
\sigma^*\mathrm{I}_2 = \mathrm{I}_1,\quad \sigma^*\mathrm{II}_2 = \mathrm{II}_1,
\label{eq:::IOJIOWAIJ2.8}
\end{align}
则存在刚体运动$\tilde{\sigma}: E^3 \to E^3$，使得$\sigma = \tilde{\sigma}\big|_{S_1}$，即曲面$S_1$和$S_2$经$E^3$的一个刚体运动后彼此重合。
\end{theorem}
\begin{remark}
该定理将唯一性推广至不同参数系的曲面情形，其证明可由读者自行完成。
\end{remark}

\begin{example}
在定理2.1的假定下，试推导由(2.2)式定义的函数$f_{\alpha\beta}(u^1,u^2)$，$f_\alpha(u^1,u^2)$，$f(u^1,u^2)$所满足的线性齐次微分方程组(2.4).
\end{example}

\begin{example}
已知函数$f_{\alpha\beta}(u^1,u^2)$，$f_\alpha(u^1,u^2)$，$f(u^1,u^2)$满足微分方程组(2.4)，命
\begin{align*}
F(u^1,u^2)=g^{\alpha\gamma}g^{\beta\delta}f_{\alpha\beta}f_{\gamma\delta}+2g^{\alpha\beta}f_\alpha f_\beta+f^2,
\end{align*}
证明：$\displaystyle\frac{\partial F(u^1,u^2)}{\partial u^\alpha}=0,\ \alpha=1,2.$
\end{example}

\begin{example}
证明定理2.2.
\end{example}



\end{document}